\section{THUẬT TOÁN NAIVE BAYES}

\subsection{Khái niệm}
Naive Bayes là một thuật toán phân lớp dựa trên định lý Bayes với giả định độc lập giữa các thuộc tính. Thuật toán này sử dụng xác suất để đoán nhãn của một mẫu mới dựa trên phân phối xác suất của các nhãn và giá trị của các thuộc tính trong tập huấn luyện.

Định lý Bayes:
$$P(A\vert{}B) = \frac{P(B\vert{}A)P(A)}{P(B)}$$

Trong đó:
\begin{itemize}
    \item $P(A\vert{}B)$: Xác suất của sự kiện $A$ xảy ra khi biết sự kiện $B$ đã xảy ra (gọi là xác suất có điều kiện).
    \item $P(B\vert{}A)$: Xác suất của sự kiện $B$ xảy ra khi biết sự kiện $A$ đã xảy ra.
    \item $P(A)$: Xác suất tiên nghiệm của sự kiện $A$ (xác suất xảy ra của $A$ mà không cần biết gì về $B$).
    \item $P(B)$: Xác suất tiên nghiệm của sự kiện $B$.
\end{itemize}

\subsection{Quy trình thực hiện}
\begin{itemize}
    \item \textbf{Bước 1:} Xác định các nhãn lớp và các thuộc tính từ tập dữ liệu huấn luyện.
    \item \textbf{Bước 2:} Tính xác suất tiên nghiệm $P(c)$ cho mỗi nhãn $c$.
    \item \textbf{Bước 3:} Tính xác suất có điều kiện $P(x\vert{}c)$ cho từng thuộc tính $x_i$ và nhãn $c$.
    \item \textbf{Bước 4:} Áp dụng công thức Bayes tính xác suất hậu nghiệm:
    $$P(c\vert{}x) = \frac{P(x\vert{}c)P(c)}{P(x)}$$
    \item \textbf{Bước 5:} Dự đoán nhãn có xác suất lớn nhất.
\end{itemize}

\subsection{Mã giả}
\begin{lstlisting}[caption={Mã giả thuật toán Naïve Bayes}]
Algorithm Naive_Bayes_Prediction:
1. Input: Training dataset, Test dataset
2. For each class C in Training dataset:
3.     Compute P(C)
4.     For each attribute x_i in X:
5.         Compute P(x_i | C)
6. For each sample X in Test dataset:
7.     Compute P(C|X) = P(C) * Product(P(x_i | C) for all x_i in X)
8.     Assign X to the class C with the highest P(C|X)
9. Output: Predicted labels
\end{lstlisting}

\subsection{Ví dụ minh họa chi tiết}

Tập huấn luyện có 14 mẫu, với 9 mẫu \textit{Yes}, 5 mẫu \textit{No}.
Xét mẫu cần phân loại gồm bốn thuộc tính:
\[X=(\mathrm{Sunny},\mathrm{Cool},\mathrm{High},\mathrm{True}).\]
\begin{table}[H]
    \centering
    \small
    \renewcommand{\arraystretch}{1.15}
    \begin{tabular}{|c|c|c|c|c|c|}
        \hline
        \rowcolor{yellow}
        \textbf{ID} & \textbf{Outlook} & \textbf{Temp} & \textbf{Humidity} & \textbf{Windy} & \textbf{Play} \\
        \hline
        1 & Sunny & Hot & High & False & No \\
        \hline
        2 & Sunny & Hot & High & True & No \\
        \hline
        3 & Overcast & Hot & High & False & Yes \\
        \hline
        4 & Rainy & Mild & High & False & Yes \\
        \hline
        5 & Rainy & Cool & Normal & False & Yes \\
        \hline
        6 & Rainy & Cool & Normal & True & No \\
        \hline
        7 & Overcast & Cool & Normal & True & Yes \\
        \hline
        8 & Sunny & Mild & High & False & No \\
        \hline
        9 & Sunny & Cool & Normal & False & Yes \\
        \hline
        10 & Rainy & Mild & Normal & False & Yes \\
        \hline
        11 & Sunny & Mild & Normal & True & Yes \\
        \hline
        12 & Overcast & Mild & High & True & Yes \\
        \hline
        13 & Overcast & Hot & Normal & False & Yes \\
        \hline
        14 & Rainy & Mild & High & True & No \\
        \hline
        \rowcolor{yellow}
        15 & Sunny & Cool & High & True & ??? \\
        \hline
    \end{tabular}
    \caption{Tập dữ liệu Play Tennis và mẫu cần phân loại}
    \label{tab:naive-bayes-play-tennis}
\end{table}
\subsubsection{Bước 1. Xác suất tiên nghiệm}
\[
P(\mathrm{Yes})=\frac{9}{14},\qquad P(\mathrm{No})=\frac{5}{14}.
\]
\subsubsection{Bước 2. Xác suất điều kiện}
\begin{center}\small
\begin{tabular}{@{}lcc@{}}\toprule
Giá trị đặc trưng&$P(\cdot\mid\mathrm{Yes})$&$P(\cdot\mid\mathrm{No})$\\\midrule
Outlook = Sunny&$2/9$&$3/5$\\
Temperature = Cool&$3/9$&$1/5$\\
Humidity = High&$3/9$&$4/5$\\
Windy = True&$3/9$&$3/5$\\\bottomrule
\end{tabular}\end{center}
\subsubsection{Bước 3. Tính xác suất hậu nghiệm}
\begin{align*}
P(\mathrm{Yes}\vert{}x)&=\frac9{14}\frac29\frac39\frac39\frac39\approx 0.00529,\\
P(\mathrm{No}\vert{}x)&=\frac5{14}\frac35\frac15\frac45\frac35\approx 0.02057.
\end{align*}
\subsubsection{Bước 4. Kết luận}
Vì $P(\mathrm{No}\vert{}x)>P(\mathrm{Yes}\vert{}x)$, mẫu được dự đoán là \textbf{No}.

\subsection{Viết chương trình cho bài toán được minh họa}
\lstinputlisting[language=Python,caption={Chương trình Python minh họa thuật toán Naive Bayes}]{source_code/naive_bayes_75.py}

\subsection{Ưu/nhược điểm của thuật toán Naive Bayes}
\begin{itemize}
    \item \textbf{Ưu điểm:}
    \begin{itemize}
        \item Hiệu quả với dữ liệu lớn và nhiều thuộc tính.
        \item Đơn giản, dễ hiểu, dễ triển khai.
        \item Thích hợp cho các bài toán phân loại văn bản như phân loại email hoặc phân tích cảm xúc.
    \end{itemize}
    \item \textbf{Nhược điểm:}
    \begin{itemize}
        \item Giả định độc lập giữa các thuộc tính thường không đúng trong thực tế, có thể làm giảm độ chính xác.
        \item Không hiệu quả khi dữ liệu có giá trị bằng 0 trong các xác suất.
    \end{itemize}
\end{itemize}